%!TEX root = ../main.tex
% APÉNDICE A — Encuesta de validación
% Fuente: docs/validacion/h3/protocolo.md del repositorio de código.

\chapter{Encuesta de validación}\label{apx:encuesta}

\begin{guia}
Los apéndices contienen material producido por el autor que respalda
el cuerpo pero lo interrumpiría. Acá va la encuesta completa: para qué
se hizo, cómo se diseñó y distribuyó, a cuántas personas llegó y el
cuestionario tal como lo vieron quienes respondieron. Los resultados
se analizan en el capítulo~\ref{cap:resultados}.
\end{guia}

\section{Introducción}

La encuesta forma parte de la validación de la hipótesis H3: la
visualización del recorrido de los datos facilita al desarrollador la
comprensión del riesgo frente a un listado de errores en texto. Se
aplica dentro de una sesión en la que cada participante resuelve dos
tareas, una con el reporte de texto de la línea de comandos y otra con
la visualización del grafo, y responde las mismas preguntas en ambas.
Incluye además una pregunta sobre la disposición a pagar el precio
definido en el capítulo~\ref{cap:economica}, que da seguimiento al
riesgo R-07.

\section{Metodología}

\paragraph{Diseño.} Estudio intra-sujeto con dos casos equivalentes y
orden contrabalanceado. El caso A contiene una inyección SQL y el caso B
una inyección de comandos; ambos tienen dos archivos, un camino
vulnerable que los cruza pasando por tres funciones y un endpoint
saneado que funciona como distractor. Cada participante recibe uno de
cuatro órdenes (tabla~\ref{tab:ordenes}), asignados en forma rotativa.

\begin{table}[H]
    \centering
    \caption{Órdenes de presentación de las condiciones y los casos}
    \label{tab:ordenes}
    \small
    \begin{tabular}{l l l}
        \toprule
        \tb{Orden} & \tb{Primera tarea} & \tb{Segunda tarea} \\
        \midrule
        O1 & Texto, caso A & Grafo, caso B \\
        O2 & Grafo, caso A & Texto, caso B \\
        O3 & Texto, caso B & Grafo, caso A \\
        O4 & Grafo, caso B & Texto, caso A \\
        \bottomrule
    \end{tabular}
\end{table}

\paragraph{Población y canal.} Personas que programan en JavaScript o
TypeScript con regularidad y que no participaron del desarrollo del
sistema. La convocatoria se realizó por mensaje directo, con sesiones
remotas de veinte minutos. \completar{Período de aplicación, cantidad de
personas convocadas y cantidad de sesiones realizadas.}

\paragraph{Registro.} Por cada tarea se registran el acierto en las
preguntas P1 a P4 (de 0 a 4, según una clave de corrección fijada antes
de las sesiones), el tiempo hasta responder P4, con un límite de seis
minutos, y las valoraciones L1 y L2.

\section{Presentación de la encuesta}

Texto leído al inicio de cada sesión:

\begin{quote}
La sesión dura veinte minutos. Vas a mirar dos fragmentos de código con
los resultados de una herramienta y responder cuatro preguntas sobre
cada uno. Se evalúa la herramienta, no a vos. Registro el tiempo y tus
respuestas sin tu nombre, y los resultados se usan solo en mi proyecto
final. Podés dejar la sesión cuando quieras. ¿Estás de acuerdo?
\end{quote}

\section{Cuestionario}

\paragraph{Preguntas de cada tarea} (abiertas; entre corchetes, el
nombre que corresponde a cada caso):
\begin{enumerate}
    \item[P1.] ¿Qué dato ingresado por el usuario llega a la operación
          peligrosa?
    \item[P2.] ¿Por qué funciones pasa ese dato, en orden, hasta llegar a
          ella?
    \item[P3.] ¿El endpoint [\texttt{verCliente} / \texttt{diagnostico}]
          tiene la misma vulnerabilidad? ¿Por qué?
    \item[P4.] Si tuvieras que corregir la vulnerabilidad con un solo
          cambio, ¿en qué archivo y línea lo harías?
\end{enumerate}

\paragraph{Valoración de cada tarea} (escala de 1, totalmente en
desacuerdo, a 5, totalmente de acuerdo):
\begin{enumerate}
    \item[L1.] Entendí por dónde viaja el dato desde la entrada hasta la
          operación peligrosa.
    \item[L2.] Tengo claro dónde aplicar la corrección.
\end{enumerate}

\paragraph{Preguntas finales:}
\begin{enumerate}
    \item[F1.] ¿Con cuál de las dos formas preferirías trabajar: texto,
          grafo o indistinto? ¿Por qué? (opción única y abierta)
    \item[F2.] Una herramienta así, que corre en el navegador sin enviar
          tu código a un servidor, ¿la pagarías a USD 8 por mes? Sí / No /
          Solo con otras funciones, ¿cuáles? (opción única y abierta)
    \item[F3.] ¿Algo que te haya costado o que agregarías? (abierta)
\end{enumerate}

\section{Casos utilizados}

\begin{lstlisting}[caption={Caso A, archivo \texttt{rutas.js}}, label={lst:caso-a-rutas}]
import { buscarPedidos, buscarCliente } from "./pedidos.js";

export function listarPedidos(req, res) {
  const estado = req.query.estado;
  const filtro = armarFiltro(estado);
  res.json(buscarPedidos(filtro));
}

export function verCliente(req, res) {
  const id = sanitize(req.params.id);
  res.json(buscarCliente(id));
}

function armarFiltro(valor) {
  return "estado = '" + valor + "'";
}
\end{lstlisting}

\begin{lstlisting}[caption={Caso A, archivo \texttt{pedidos.js}}, label={lst:caso-a-pedidos}]
export function buscarPedidos(filtro) {
  return db.query("SELECT * FROM pedidos WHERE " + filtro);
}

export function buscarCliente(id) {
  return db.query("SELECT * FROM clientes WHERE id = " + id);
}
\end{lstlisting}

El caso B tiene la misma estructura: \texttt{exportarReporte} recibe
\texttt{req.body.carpeta}, la pasa por \texttt{armarRuta} y llega a
\texttt{execSync} en \texttt{comprimirCarpeta}; el distractor
\texttt{diagnostico} pasa el host por \texttt{validate}. Los archivos de
ambos casos y los reportes de texto entregados a los participantes se
encuentran en la carpeta \texttt{docs/validacion/h3/} del repositorio
(anexo~\ref{anx:repositorios}).
